\documentclass[11pt,a4paper]{article}
\usepackage[margin=26mm]{geometry}
\usepackage{amsmath,amssymb,amsthm,microtype}
\usepackage[hidelinks]{hyperref}
\newcommand{\E}{\mathbb E}\newcommand{\Pp}{\mathbb P}
\newcommand{\CE}{C_{\rm E}}
\newtheorem{theorem}{Theorem}\newtheorem{lemma}[theorem]{Lemma}
\title{The sharp Berry--Esseen bound survives small independent noise}
\author{}\date{28 October 2026}
\begin{document}\maketitle

If the two conditional noise laws are identical, the marginal is a
Bernoulli variable plus independent noise. Conditioning on that noise
allows us to apply the sharp binomial theorem at a random threshold.
A finite Taylor replacement then compares the smoothed Gaussian term
with the Gaussian of the correct total variance. The increase in the
third absolute moment pays for this error. The resulting argument covers
every sample size and threshold in this subclass.

\begin{theorem}\label{thm:closure}
Let $B\sim\mathrm{Bern}(p)$, where $1/4\le p\le3/4$, and let $U$
be independent of $B$, centered, and supported on an interval of length
at most $1/5$. Put
\[
 v=p(1-p),\quad s=\E U^2,\quad
 X=\frac{B-p+U}{\sqrt{v+s}},\quad \rho=\E|X|^3,\qquad
 \CE=\frac{3+\sqrt{10}}{6\sqrt{2\pi}}.
\]
For independent copies $X_1,\ldots,X_n$, every positive integer $n$ satisfies
\begin{equation}\label{eq:main}
 \sqrt n\,\sup_{H\in\mathbb R}
 \left|\Pp(X_1+\cdots+X_n\le H)-\Phi(H/\sqrt n)\right|
 \le \CE\rho-\frac{3s}{1000v^{3/2}}.
\end{equation}
In particular the sharp Berry--Esseen bound holds for this class, with
a strict relative-noise margin when $s>0$. No upper bound on $ns$,
lower bound on $n$, threshold restriction, acquisition or contact
hypothesis is required.
\end{theorem}

\section{The two estimates}
The only external probability theorem used is Schulz's sharp binomial
bound \cite[Theorem 1, printed p.~1]{Schulz}. For $B_n\sim\mathrm{Bin}(n,p)$,
$0<p<1$ and every $n\ge1$, it states
\begin{equation}\label{eq:schulz}
 \sup_{z\in\mathbb R}
 \left|\Pp(B_n\le z)-\Phi\!\left(\frac{z-np}{\sqrt{nv}}\right)\right|
 <\frac{\CE\rho_B}{\sqrt n},
 \qquad \rho_B=\frac{1-2v}{\sqrt v}.
\end{equation}
This uses closed distribution functions and includes both signs and all
real thresholds. No restricted-parameter sharpening is invoked.

\begin{lemma}[A smooth replacement with its complete remainder]
\label{lem:smooth}
Let $U_1,\ldots,U_n$ be independent copies of a centered variable with
variance $s$ and finite third absolute moment $m_3$.
For every real $t$,
\begin{align}\label{eq:smooth}
 &\left|\E\Phi\!\left(\frac{t-np-\sum_iU_i}{\sqrt{nv}}\right)
       -\Phi\!\left(\frac{t-np}{\sqrt{n(v+s)}}\right)\right|\\
 &\hspace{25mm}\le
 \frac{\phi(0)}{6\sqrt n\,v^{3/2}}
                  \left(m_3+\sqrt{\frac8\pi}\,s^{3/2}\right).\notag
\end{align}
\end{lemma}
\begin{proof}
For $f(w)=\Phi((t-np-w)/\sqrt{nv})$,
\[
       \|f'''\|_\infty\le\frac{\phi(0)}{(nv)^{3/2}}.
\]
Indeed $\Phi'''(x)=(x^2-1)\phi(x)$ has absolute maximum $\phi(0)$:
the other possible nonzero extrema are at $x=\pm\sqrt3$, with value
$2e^{-3/2}\phi(0)<\phi(0)$.

Replace $U_1,\ldots,U_n$ one at a time by independent
$G_i\sim N(0,s)$. Conditional on all the other terms, Taylor's formula
through degree two has remainder bounded by
$\|f'''\|_\infty|x|^3/6$. Means and second moments cancel.
Each replacement therefore costs at most
$\|f'''\|_\infty(m_3+\E|G_i|^3)/6$.
The Gaussian third absolute moment is
$\E|G_i|^3=\sqrt{8/\pi}\,s^{3/2}$. Summing the $n$ remainders gives
the stated bound, since a Gaussian sum of variance $ns$ convolved
with the Gaussian CDF of variance $nv$ gives variance $n(v+s)$.
The argument is finite, uniform in $t$, and remains valid at $s=0$.
\end{proof}

\begin{lemma}[The moment increase pays for smoothing]\label{lem:moment}
Under Theorem~\ref{thm:closure}'s hypotheses,
\begin{equation}\label{eq:moment}
 \rho-\rho_B\ge\frac{27}{400}\frac{s}{v^{3/2}},
 \qquad
 \frac{\phi(0)}6\left(m_3+\sqrt{\frac8\pi}s^{3/2}\right)
       \le\frac3{125}s .
\end{equation}
\end{lemma}
\begin{proof}
Put $m_3=\E|U|^3$ and $\mu_3=\E U^3$. Write the containing interval as
$[-a,b]$, with $a,b\ge0$ and $d=a+b\le1/5$; centering permits this choice. Then
\[
 |U|\le d,\qquad s\le ab\le d^2/4\le1/100,\qquad
 m_3\le ds,\qquad |\mu_3|\le ds.
\]
The variance bound follows by averaging
$(U+a)(b-U)\ge0$.

Since $|U|\le1/5<\min(p,1-p)$, the sign of $B-p+U$ is fixed by $B$.
Expanding the two cubics gives the exact identity
\begin{equation}\label{eq:rawmoment}
 \E|B-p+U|^3=v(1-2v)+6vs+(2p-1)\mu_3.
\end{equation}
For $x\ge0$, Taylor's formula gives
$(1+x)^{3/2}\le1+\frac32x+\frac38x^2$, because the third derivative
is negative. Consequently
\begin{align*}
 (v+s)^{3/2}(\rho-\rho_B)
 &\ge s\left[9v-\frac32-|2p-1|d
              -\frac38(1-2v)\frac{s}{v}\right]\\
 &\ge s\left[\frac3{16}-\frac1{10}-\frac1{80}\right]
   =\frac3{40}s .
\end{align*}
Here $v\ge3/16$, $|2p-1|\le1/2$, $1-2v\le5/8$ and $s/v\le4/75$.
Also
\[
 \left(\frac{v}{v+s}\right)^{3/2}
 \ge\left(\frac{75}{79}\right)^{3/2}>\frac9{10}.
\]
This proves the first bound in \eqref{eq:moment}.

Finally $\phi(0)<2/5$, $\sqrt{8/\pi}<8/5$, and $\sqrt s\le1/10$ give
\[
 \frac{\phi(0)}6
       \left(m_3+\sqrt{\frac8\pi}s^{3/2}\right)
 \le\frac1{15}\left(\frac15+\frac85\frac1{10}\right)s
 =\frac3{125}s.
\]
All constants are independent of $n$, the threshold and the law.
\end{proof}

\section{Completion and scope}
\begin{proof}[Proof of Theorem~\ref{thm:closure}]
Let $Y_i=B_i+U_i$, with all Bernoulli and noise variables independent.
For every raw threshold $t$, condition on $W=\sum_iU_i$ and use
\eqref{eq:schulz}. Lemma~\ref{lem:smooth} then gives
\[
 \sqrt n\left|
 \Pp\Big(\sum_iY_i\le t\Big)
       -\Phi\!\left(\frac{t-np}{\sqrt{n(v+s)}}\right)\right|
 \le\CE\rho_B+\frac3{125}\frac{s}{v^{3/2}}.
\]
The elementary inequality $\CE>2/5$ and Lemma~\ref{lem:moment} imply
\[
 \CE(\rho-\rho_B)\ge\frac{27}{1000}\frac{s}{v^{3/2}}.
\]
Since $27/1000-3/125=3/1000$, subtraction proves \eqref{eq:main}
after setting $t=np+H\sqrt{v+s}$ and taking the supremum.
Conditioning retains every threshold atom; only the smooth Gaussian
term was differentiated.
\end{proof}

The constant estimates need only $25/8<\pi<22/7$. The upper bound follows
from
\[
 22/7-\pi=\int_0^1\frac{x^4(1-x)^4}{1+x^2}\,dx>0;
\]
the integral of the alternating geometric polynomial through $-x^{126}$
in $(1+x^2)^{-1}$ gives the lower bound. Moreover
$\CE^2>37/(72\pi)>259/1584>4/25$, since $\sqrt{10}>3$.
The companion rational checker verifies these finite inequalities and
the exact cubic expansions on actual finite laws. It does not reprove
Schulz's theorem.

If $U_0$ and $U_1$ are the two conditional noise laws in the acquired
Bernoulli coordinates and their laws coincide, then $B$ and $U$ are
independent. The common support lies in $[-a_1,a_0]$, so its width is
at most $.072$ on the right and $.068$ on the left. The theorem therefore
excludes every such acquired contrary witness, including positive noise
above the currently certified regional ceilings. It also applies well
outside that acquired setting.

This removes the count and Poisson comparison computations for precisely
this common-noise subclass. It does not show that every extremizer has
equal conditional noise laws, nor does it prove the finite base or the
global conjecture. It bounds the sharp Berry--Esseen excess; it does not
assert domination by the discrepancy at every chosen Bernoulli anchor.
The unweighted crossing estimate and its phase counterexamples in
\texttt{loss-common.tex} are supplementary and are not dependencies of
this proof.

\begin{thebibliography}{1}
\bibitem{Schulz}
J. Schulz, \emph{The optimal Berry--Esseen constant in the binomial case},
doctoral dissertation, Universit\"at Trier, 2016.
Theorem 1, printed p.~1 (PDF page 11).
\url{https://d-nb.info/1197702695/34}.
\end{thebibliography}
\end{document}
